%%FRAG s:ewa
%%BEFORE: \subsubsection{Khai báo đầu vào, uniform và biến ra}
\subsubsection{Cơ sở toán: EWA Splatting, chiếu ellipsoid 3D thành ellipse 2D}\label{s:ewa}
Mục này giải thích toán của pass 1 (\code{splat_prep.vert}); phần chiếu và trộn nằm ở \code{splat.vert}, \code{splat.frag} (mục \ref{s:blend}).

\paragraph{Hiệp phương sai 3D.}
$\Sig = R S S^{\top} R^{\top} = M M^{\top}$, với $M=R\,\mathrm{diag}(s)$ (shader: \code{M = R * mat3(scale...)}, \code{sigma = M * transpose(M)}). $\Sig$ đối xứng bán xác định dương; trục ellipsoid là cột của $R$, bán trục là $s$.

\paragraph{Jacobian của phép chiếu.}
Điểm ảnh: $x_{px}=f_{px}v_x/(-v_z)$, $y_{px}=f_{px}v_y/(-v_z)$. Đạo hàm theo $(v_x,v_y,v_z)$ tại tâm:
\[
J=\begin{bmatrix}\frac{f}{-v_z}&0&\frac{f v_x}{v_z^2}\\0&\frac{f}{-v_z}&\frac{f v_y}{v_z^2}\\0&0&0\end{bmatrix},\qquad
f_{px}=\tfrac H2 P_{11}=\frac{H/2}{\tan(\mathrm{fov}/2)}
\]
(\code{splat_renderer}: \code{uFocal = height/2 * projection[1,1]}). Phép chia cho $c_w$ là phi tuyến, nên dùng \textbf{xấp xỉ tuyến tính bậc nhất} tại tâm (khai triển Taylor) để ellipsoid vẫn cho ra ellipse.

\paragraph{Hiệp phương sai 2D.}
\[
\Sig' = J\,W\,\Sig\,W^{\top}J^{\top}.
\]
Trong shader gộp thêm ma trận layer $sA$: \code{T = J * mat3(ModelView) * linear}, \code{cov = T Σ Tᵀ} ($\Sig$ ở tọa độ file; \code{linear = s·A} đã gồm co giãn nên không cần $s^2$ riêng). Lấy khối $2\times2$: $a=\mathrm{cov}_{00}$, $b=\mathrm{cov}_{01}$, $c=\mathrm{cov}_{11}$.

\textbf{Bộ lọc thông thấp:} $a\mathrel{+}=0.3$, $c\mathrel{+}=0.3$ (tương đương cộng $0.3I$), bảo đảm gaussian rộng ít nhất khoảng 1 pixel, chống răng cưa như EWA gốc và 3DGS.

\paragraph{Trị riêng và vector riêng của ma trận 2x2.}
\begin{align*}
\mathrm{mid}&=\tfrac{a+c}{2}, & \mathrm{radius}&=\sqrt{\left(\tfrac{a-c}{2}\right)^2+b^2},\
\lambda_1&=\mathrm{mid}+\mathrm{radius}, & \lambda_2&=\max(\mathrm{mid}-\mathrm{radius},\,0.1)\ \text{(kẹp chống suy biến)},\
\mathrm{axis}_1&=\mathrm{normalize}(b,\lambda_1-a), & \mathrm{axis}_2&=(-\mathrm{axis}_{1,y},\mathrm{axis}_{1,x}).
\end{align*}
Nếu $|b|<10^{-6}$ thì dùng trục $x$ hoặc $y$. Kiểm tra: $(A-\lambda_1 I)v=0\Rightarrow(a-\lambda_1)v_1+bv_2=0\Rightarrow v\propto(b,\lambda_1-a)$. Bán trục ellipse là $\sqrt{\lambda_1},\sqrt{\lambda_2}$ (độ lệch chuẩn theo hai hướng riêng).

\paragraph{Kích thước hình vuông (billboard).}
\begin{itemize}
\item $\mathrm{extent}=\min\!\big(\sqrt{2\ln(255\alpha)},\,3\big)$, từ phương trình $\alpha e^{-r^2/2}=1/255\Rightarrow r^2=2\ln(255\alpha)$. Gaussian mờ cho hình vuông nhỏ; tối đa $3\sigma$.
\item $\mathrm{major}=\min(\mathrm{extent}\cdot\sqrt{\lambda_1},1024)\cdot\mathrm{axis}_1\cdot(2/\mathrm{viewport})$, \code{minor} tương tự với $\lambda_2$. Nhân $2/W$, $2/H$ để đổi pixel sang NDC (NDC rộng 2).
\item \textbf{Loại bỏ (cull):} $\alpha<1/255$; $c_w\le0$ (sau camera); $|\mathrm{ndc}_x|$ hoặc $|\mathrm{ndc}_y|>1.3$ (đệm 30\% vì ellipse có thể lấn vào màn hình).
\end{itemize}
%%END

%%FRAG s:blend
%%BEFORE: \subsection{\texorpdfstring{\code{splat.frag}}{splat.frag}}
\subsubsection{Cơ sở toán: hiển thị gaussian (pass 2, đỉnh)}\label{s:blend}
Mỗi instance vẽ 4 góc $(\pm1,\pm1)$: \code{pos_ndc = center + cx*major + cy*minor}; \code{vCorner = (cx,cy)*extent} (đơn vị $\sigma$ trong hệ trục ellipse). Fragment shader và phép trộn alpha xem tiếp ở \code{splat.frag} ngay sau đây; kích thước \code{major}, \code{minor}, \code{extent} lấy từ mục \ref{s:ewa}.
%%END

%%FRAG s:blend-frag
%%BEFORE: \subsection{\texorpdfstring{\code{shadow.vert}}{shadow.vert}}
\subsubsection{Cơ sở toán: fragment và trộn alpha ``over''}\label{s:blend-frag}
\paragraph{Fragment (\code{splat.frag}).}
\begin{itemize}
\item $r^2=|\mathrm{vCorner}|^2$; bỏ nếu $r^2>9$ (ngoài $3\sigma$).
\item $a=\alpha\exp(-r^2/2)$ (mật độ Gauss 2D chuẩn hóa về đỉnh bằng $\alpha$). Bỏ nếu $a<1/255$.
\item Xuất $(\mathrm{rgb}\cdot a,\,a)$, tức \textbf{premultiplied alpha}.
\end{itemize}

\paragraph{Trộn ``over'' (\code{glBlendFunc(GL_ONE, GL_ONE_MINUS_SRC_ALPHA)}).}
\[
C_i = c_i a_i + (1-a_i)\,C_{i-1}
\]
(với rgb đã nhân $a$: $C=\mathrm{src.rgb}+(1-\mathrm{src.a})\,\mathrm{dst}$). Phép này không giao hoán, nên phải vẽ \textbf{từ xa đến gần}. Depth test tắt (vật trong suốt không được che nhau).
%%END

%%FRAG s:viewport
%%BEFORE: \subsubsection*{Khai báo mức module (dòng 21--28)}
\subsubsection{Cơ sở toán: biến đổi viewport (\texorpdfstring{\code{brush._to_screen}}{brush._to_screen})}\label{s:viewport}
\[
x_{px}=x_0+\tfrac{n_x+1}{2}W,\qquad y_{px}=y_0+\tfrac{1-n_y}{2}H\ \text{(đảo $y$: màn hình có $y$ hướng xuống)}.
\]
%%END

%%FRAG s:budget
%%BEFORE: \subsubsection*{Khai báo mức module (dòng 35--54)}
\subsubsection{Cơ sở toán: tự chỉnh ngân sách (\texorpdfstring{\code{SplatRenderer.adapt_budget}}{SplatRenderer.adapt_budget})}\label{s:budget}
Đo thời gian GPU bằng query. $\mathrm{cost}=\mathrm{median}(\mathrm{ms}/\mathrm{count})$; $\mathrm{wanted}=\mathrm{target\_ms}/\mathrm{cost}$. Nếu lệch dưới 10\% thì giữ nguyên; ngược lại $\mathrm{step}=\mathrm{clip}(\mathrm{wanted}/\mathrm{current},0.75,1.33)^{1/2}$ (căn bậc hai nên đổi tối đa khoảng 15\%/lần, làm trơn để tránh dao động); $\mathrm{budget}=\mathrm{clip}(\mathrm{current}\cdot\mathrm{step},\,\mathrm{MIN\_BUDGET},\,\mathrm{count})$.
%%END
